\documentclass[11pt]{letter}
\usepackage[margin=1in]{geometry}
\usepackage{times}

\signature{Daniel Polónia, Ph.D. Candidate}
\address{
    Department of Economics, Management, Industrial Engineering and Tourism \\
    Universidade de Aveiro \\
    Campus Universitário de Santiago \\
    3810-193 Aveiro, Portugal \\
    Email: dpolonia@ua.pt
}

\begin{document}

\begin{letter}{
Editor-in-Chief \\
Health Care Management Science \\
Springer Nature
}

\opening{Dear Editor,}

I am pleased to submit for consideration the manuscript titled \textbf{``Stakeholder-Distributed Distress: Measuring Financial Sustainability in Public Hospitals Under Soft Budget Constraints''} for publication in \textit{Health Care Management Science}.

This study makes three novel contributions at the intersection of corporate finance, health economics, and management science. First, we develop and validate the Public Hospital Financial Sustainability Index (PHFSI), a theoretically grounded composite measure specifically designed for organizations operating under soft budget constraints where conventional bankruptcy prediction models fail. Second, using an eight-year panel of 741 hospital-year observations from Portugal's National Health Service (2017--2024), we provide rigorous econometric evidence of subsidy-induced moral hazard in healthcare finance: a one-percentage-point increase in subsidy dependence reduces financial sustainability by 0.547 standard deviations (p < 0.001). Third, we introduce a stakeholder-distributed distress framework formalizing how public hospital financial stress transfers sequentially across suppliers, staff, patients, and taxpayers---fundamentally different from private sector bankruptcy where costs concentrate at discrete resolution events.

This manuscript is particularly timely and policy-relevant for three reasons:

\textbf{1. Global Relevance.} Public hospitals in Beveridgean healthcare systems worldwide---the UK's NHS, Spain's regional health services, Italy's SSN, Greece's ESY, and Canadian provincial hospitals---exhibit soft budget constraint dynamics. The €500 million October 2024 capital injection into Portuguese SNS hospitals mirrors recent bailouts in Italy (€3.2 billion, 2023), Spain (€2.8 billion, 2022), and Greece (€1.1 billion, 2021). Our framework and measurement approach generalize across these contexts.

\textbf{2. Methodological Innovation.} Traditional bankruptcy prediction models (Altman's Z-score, Ohlson's O-score, machine learning approaches) achieve areas under the ROC curve near 0.50 (random guessing) for public hospitals because they assume distress culminates in bankruptcy---an institutionally impossible event in mission-critical services. The PHFSI reconceptualizes distress measurement by explicitly incorporating stakeholder burden indicators (supplier payment delays, staff retention stress, patient quality deterioration) alongside conventional balance sheet ratios.

\textbf{3. Policy Impact.} Our findings have direct implications for healthcare financial management and policy reform. We estimate that implementing PHFSI-based early warning systems could reduce late-stage bailout costs by 15--20\% through preemptive intervention. The strong subsidy-payment delay relationship (β = 87.3 days per percentage point subsidy increase, p < 0.001) suggests that payment mechanism reforms---accelerating government reimbursements and introducing subsidy conditionality---could reduce moral hazard without compromising access.

The manuscript fits squarely within \textit{Health Care Management Science}'s scope of publishing ``high-quality research that contributes to the understanding and improvement of the management and delivery of health care services.'' Our work addresses critical questions in healthcare financial management, uses rigorous econometric methods (two-way fixed effects panel models with robustness checks), and provides actionable policy recommendations for health system administrators and policymakers.

Methodologically, the study employs publicly available data from Portugal's SNS Transparency Portal, Eurostat, and national statistics---ensuring full replication potential. All analyses use robust standard errors clustered at the hospital level, and results are validated across eight robustness specifications. The PHFSI's modular design allows adaptation to different institutional contexts while preserving theoretical coherence.

This manuscript has not been published previously and is not under consideration elsewhere. All authors have approved the manuscript and agree with its submission to \textit{Health Care Management Science}. The research complies with ethical standards and uses only publicly available administrative data requiring no institutional review board approval.

For your convenience, I suggest the following potential reviewers with expertise in healthcare finance, soft budget constraints, and distress prediction:

\begin{enumerate}
    \item \textbf{Professor Pedro Pita Barros}, Nova School of Business and Economics, Portugal (pedro.barros@novasbe.pt) --- Expert in healthcare economics and Portuguese SNS institutional context.

    \item \textbf{Professor Martin Gaynor}, Carnegie Mellon University, USA (mgaynor@cmu.edu) --- Leading scholar on hospital competition and healthcare markets.

    \item \textbf{Professor Carol Propper}, Imperial College Business School, UK (c.propper@imperial.ac.uk) --- Expert on NHS financial performance and quality under budget constraints.

    \item \textbf{Professor Janos Kornai}, Harvard University (emeritus), USA (kornai@harvard.edu) --- Originator of soft budget constraint theory.

    \item \textbf{Professor Edward Altman}, New York University Stern School of Business, USA (ealtman@stern.nyu.edu) --- Developer of corporate distress prediction models.
\end{enumerate}

I am happy to provide any additional information or clarification you may need. Thank you for considering this manuscript for publication in \textit{Health Care Management Science}.

\closing{Sincerely,}

\end{letter}

\end{document}
